\section{Mathematical Setup}
\label{sec:setup}

\subsection{QGPV equation on the \texorpdfstring{$\beta$}{beta}-plane}

The barotropic quasi-geostrophic potential vorticity (QGPV) equation on the
$\beta$-plane reads
\begin{equation}\label{eq:qgpv}
  \frac{Dq}{Dt} = 0, \qquad q = \nabla^2\psi + \beta y,
\end{equation}
where $\psi$ is the stream function, $\beta = \dd f/\dd y$ is the
meridional gradient of the Coriolis parameter, and $D/Dt = \partial_t +
J(\psi, \cdot)$ is the material derivative.

\subsection{Log-polar coordinate transformation}

Under the conformal map $(x,y) \mapsto (\rho,\theta)$ defined by
$x + iy = e^{\rho + i\theta}$, the Laplacian transforms as
\begin{equation}\label{eq:laplacian-logpolar}
  \nabla^2 = e^{-2\rho}\!\left(
    \frac{\partial^2}{\partial\rho^2} +
    \frac{\partial^2}{\partial\theta^2}
  \right).
\end{equation}
Linearizing \eqref{eq:qgpv} about a zonal base state $U(\rho)$ and Fourier
expanding the perturbation $\psi' = \psihat_n(\rho)\,e^{in\theta}\,
e^{-i\omega t}$ yields the Rayleigh--Kuo eigenvalue ODE
\begin{equation}\label{eq:rayleigh-kuo}
  (U - c)\bigl(\psihat_n'' - n^2 \psihat_n\bigr)
  + (\beta - U'')\,\psihat_n = 0,
\end{equation}
where $c = \omega/n$ is the phase speed and primes denote $\dd/\dd\rho$.
For stationary waves ($c = 0$) this simplifies to
\begin{equation}\label{eq:stationary-ode}
  \psihat_n'' + \left[\frac{\beta - U''}{U} - n^2\right] \psihat_n = 0.
\end{equation}

\subsection{M\"obius transformations and the conjugacy theorem}
\label{sec:mobius-conjugacy}

A M\"obius (linear fractional) transformation
$f(z) = (az+b)/(cz+d)$ with $ad - bc \neq 0$ acts on the Riemann sphere
$\mathbb{C}_\infty = \mathbb{C} \cup \{\infty\}$.  Its classification is
determined by the \emph{multiplier} $\lam$, defined as follows.

\begin{theorem}[M\"obius conjugacy]\label{thm:conjugacy}
Let $f$ be a non-parabolic M\"obius transformation with distinct fixed
points $z_1, z_2 \in \mathbb{C}_\infty$.  Define the conjugating map
$g(z) = (z - z_1)/(z - z_2)$.  Then:
\begin{enumerate}
  \item $h = g \circ f \circ g^{-1}$ satisfies $h(w) = \lam w$, where
    $\lam = f'(z_1)$ is the derivative at the fixed point $z_1$.
  \item The multiplier satisfies $\lam = \mu_1/\mu_2$ where $\mu_1, \mu_2$
    are the eigenvalues of the associated matrix
    $\bigl[\begin{smallmatrix}a&b\\c&d\end{smallmatrix}\bigr]$.
  \item Writing $\lam = e^{\sigma + i\alpha}$ with
    $\sigma = \log|\lam|$, $\alpha = \arg(\lam)$:
    \begin{center}
    \begin{tabular}{lll}
    \toprule
    $\sigma \neq 0$, $\alpha \neq 0$ & loxodromic & logarithmic spiral orbits \\
    $\sigma = 0$, $\alpha \neq 0$    & elliptic   & circular orbits \\
    $\sigma \neq 0$, $\alpha = 0$    & hyperbolic & radial orbits \\
    \bottomrule
    \end{tabular}
    \end{center}
\end{enumerate}
\end{theorem}

\begin{proof}
For $c \neq 0$, the fixed points of $f$ satisfy $cz^2 + (d-a)z - b = 0$
with discriminant $\Delta = (a+d)^2 - 4(ad-bc) \neq 0$ (non-parabolic),
giving two distinct finite fixed points.  For $c = 0$, $f$ is affine with
one finite fixed point and $z_2 = \infty$; the conjugacy proceeds
identically with $g(z) = z - z_1$.
Since $g(z_1) = 0$ and $g(z_2) = \infty$, the map $h$ fixes both $0$ and
$\infty$.  The only M\"obius maps fixing both are $w \mapsto \lam w$.  By
the chain rule, $\lam = h'(0) = g'(z_1) \cdot f'(z_1) \cdot
(g^{-1})'(0) = f'(z_1)$.
\end{proof}

\begin{remark}[Fixed point character]\label{rem:fixed-points}
For loxodromic and hyperbolic transformations ($|\lam| \neq 1$), $z_1$
is repelling and $z_2$ is attracting (or vice versa, depending on the
convention $|\lam| \ge 1$).  For \emph{elliptic} transformations
($|\lam| = 1$), both fixed points are \emph{neutral}---centers of
rotation, neither attracting nor repelling.  The term ``repelling fixed
point'' used in the loxodromic context does not apply to the elliptic
case.
\end{remark}

\subsection{From circles to polygons: the \texorpdfstring{$\mathbb{Z}_N$}{Z\_N}
projection}
\label{sec:circle-to-polygon}

The conjugacy theorem establishes that elliptic ($\sigma = 0$) orbits are
\emph{circles}, not polygons.  The regular $N$-gon arises through a
separate geometric step.

\begin{proposition}[Polygonal meander from $\mathbb{Z}_N$ perturbation]
\label{prop:polygon}
Any axisymmetric jet perturbed by an $n = N$ Fourier mode produces a
jet boundary
\begin{equation}\label{eq:jet-boundary}
  \rho_{\mathrm{boundary}}(\theta) = \rhostar + \varepsilon\cos(N\theta + \phi)
\end{equation}
that is a meander with $N$-fold symmetry, approximating a regular $N$-gon
for small $\varepsilon$.
\end{proposition}

\begin{proof}
The total PV field with the $n = N$ Fourier perturbation is
$q(\rho, \theta) = q_0(\rho) + q_n(\rho)\cos(N\theta + \phi)$.
The jet boundary is the PV contour $q = q_B$ (constant).  Linearizing
around $\rho = \rhostar$ where $q_0(\rhostar) = q_B$:
\[
  q_0'(\rhostar)\,\delta\rho + q_n(\rhostar)\cos(N\theta + \phi) = 0,
\]
giving $\delta\rho(\theta) = -\frac{q_n(\rhostar)}{q_0'(\rhostar)}
\cos(N\theta + \phi)$, which is \eqref{eq:jet-boundary} with
$\varepsilon = |q_n(\rhostar)/q_0'(\rhostar)|$.

This meander has $N$ maxima at $\theta = 2\pi k/N$ ($k = 0, \ldots, N-1$),
all at the same radius $R(1 + \varepsilon)$, equally spaced in angle---the
vertices of a regular $N$-gon to leading order in $\varepsilon$.  (At
$O(\varepsilon^2)$, the chord lengths between adjacent vertices differ
because the meander curves between maxima; the polygon approximation is
valid for $\varepsilon \ll 1$.)
\end{proof}

\begin{remark}[Role of Proposition~\ref{prop:polygon}]
This result is a standard fact about Fourier perturbations and does not
require M\"obius theory.  Its role in the paper is to provide the
geometric bridge between the M\"obius conjugacy theorem (which gives
\emph{circles} at $\sigma = 0$) and the observed \emph{polygons}: the
polygon arises when the QGPV selects a specific azimuthal mode $n = N$
(\cref{sec:theorems}), which converts the circular base state into an
$N$-gon meander via this proposition.
\end{remark}

This distinction matters: the M\"obius orbit is a circle; the polygon
emerges from the \emph{interaction} between the circular orbit structure
and the discrete $\mathbb{Z}_N$ symmetry imposed by the wavenumber
selection (Rossby dispersion or vortex dynamics).  The two steps are:
\begin{enumerate}
  \item $\sigma = 0$ places the system on the circular-orbit family
    (M\"obius conjugacy theorem);
  \item $\alpha = 2\pi/N$ restricts this to the $\mathbb{Z}_N$-symmetric
    sector, producing the $N$-gon meander.
\end{enumerate}

\subsection{Loxodromic flow family and orbit structure}

The continuous one-parameter group generated by $h(w) = \lam w$ is
$w(t) = e^{(\sigma + i\alpha)t}\,w_0$.  In polar coordinates
$w = Re^{i\Theta}$:
\begin{equation}\label{eq:orbit-polar}
  R(t) = R_0\,e^{\sigma t}, \qquad \Theta(t) = \alpha t + \Theta_0.
\end{equation}
Eliminating $t$ (for $\alpha \neq 0$):
\begin{equation}\label{eq:spiral-equation}
  \ln R = \ln R_0 + \frac{\sigma}{\alpha}(\Theta - \Theta_0).
\end{equation}
This is a logarithmic spiral with constant winding angle
$\psi = \arctan(\alpha/\sigma)$.  At $\sigma = 0$: $R = R_0$ for all
$\Theta$ (circular orbit, $\psi = \pi/2$).  As $\sigma \to 0^+$: the
spiral opens continuously to a circle.  Both $R(\Theta; \sigma)$ and
$\psi(\sigma)$ are $C^\infty$ in $\sigma$ for fixed $\alpha \neq 0$.

\begin{remark}[Stream function and the periodic domain]
\label{rem:stream-function}
The full complex potential $w(z) = Az^s$ with $s = \sigma + i\alpha$ gives
\begin{equation}\label{eq:psi-full}
  \psi_{\mathrm{full}}(\rho,\theta)
  = |A|\,e^{\sigma\rho - \alpha\theta}
  \sin(\alpha\rho + \sigma\theta + \arg A).
\end{equation}
The factor $e^{-\alpha\theta}$ is not $2\pi$-periodic in $\theta$ for
$\alpha \neq 0$.  This function is properly defined on the
\emph{universal cover} of the punctured plane---the strip
$\{(\rho,\theta) : \rho \in \mathbb{R},\; \theta \in \mathbb{R}\}$
with the identification $z = e^{\rho + i\theta}$---rather than on
the periodic domain $\theta \in [0, 2\pi)$.

\emph{We do not use $\psi_{\mathrm{full}}$ as a periodic stream function.}
The QGPV analysis proceeds entirely through the Fourier-decomposed
eigenvalue equation \eqref{eq:rayleigh-kuo}, which governs the radial
eigenfunction $\psihat_n(\rho)$ for each azimuthal mode $n$ independently.
The full perturbation
\begin{equation}\label{eq:fourier-mode}
  \psi'(\rho,\theta,t) = \psihat_n(\rho)\,e^{in\theta}\,e^{-i\omega t}
\end{equation}
\emph{is} $2\pi$-periodic in $\theta$ and satisfies the linearized QGPV
by construction (it is the Fourier decomposition of the linearized
equation).

The complex potential $w(z) = Az^s$ enters only through the \emph{orbit
geometry} of the M\"obius transformation $f(z) = \lam z$
(\cref{sec:mobius-conjugacy})---specifically, the shape of the orbits
as a function of $\sigma$---and not as a stream function substituted
into the QGPV.  The bridge between the M\"obius geometry and the QGPV
is the stationarity theorem (Theorem~\ref{thm:stationarity}), which
operates on $\psihat_n(\rho)$, not on $\psi_{\mathrm{full}}$.
\end{remark}
